% Section IV -- Results. HAR is the primary dataset; CIFAR-10 is secondary.
%
% PROVENANCE OF EVERY NUMBER IN THIS FILE
%   HAR per-round energy (Table~\ref{tab:perround}): DERIVED here from the
%     run totals and the median idle baseline reported by the Phase-1 check
%     on the 18-run HAR sweep. Replace with `python analyze.py results_har/`
%     output, which subtracts each run's OWN idle baseline rather than the
%     median and will shift the last digit.
%   HAR stopping rounds for eps in {inf, 1}: single-seed pilots at 100 rounds.
%     The sweep's own three-seed table supersedes them -- see the TODO below.
%   CIFAR-10 numbers: `python analyze.py results/` on the 18-run CIFAR sweep.
%   Ablations: CIFAR only so far; the HAR ablation grid has not been run.
%
% ------------------------------------------------------------------ TODO
% Run and paste in:
%   python analyze.py results_har/
%   python analyze.py results_har/ --targets 0.40,0.50,0.85
% Then fill: Table~\ref{tab:stopping} (all six epsilon, three seeds),
% Table~\ref{tab:target}, Table~\ref{tab:carbon}, and replace the derived
% figures in Table~\ref{tab:perround}. Search this file for TODO.
% -----------------------------------------------------------------------

\section{Results}

We report on two datasets. UCI-HAR is primary: clients are the individual
volunteers the data was collected from, so the client boundary is a real
person rather than a synthetic partition, and the task is the sensor-on-edge
setting the study is about. CIFAR-10 is secondary and serves as a harder
benchmark on which the same effects can be checked.

All figures are means over three seeds unless stated. The measurement noise
floor, computed from the across-seed spread within each condition, is
$\pm 5.43$~W ($9.46\%$) on HAR and $\pm 4.38$~W ($6.38\%$) on CIFAR-10; no
difference smaller than that is treated as resolvable. HAR rounds are short
(roughly $0.48$~s), so per-round energy there is a noisier estimate than on
CIFAR-10 and whole-run totals are the more reliable quantity.

\subsection{The Per-Round Cost of Privacy Is Fixed}

Table~\ref{tab:perround} gives energy per round on HAR. Net of the idle
baseline, private training costs $21.1$~J per round on average against
$14.6$~J without privacy, an overhead of $+45\%$. Across the five private
conditions the per-round figure varies by $10.1\%$ of its mean against a
$9.46\%$ noise floor --- within resolution, though only just, and this is the
one headline number the shorter HAR rounds make least comfortable. CIFAR-10
gives the same result with more margin: $+60\%$ overhead, $9.6\%$ spread
against a $6.38\%$ floor.

Wall-clock says the same thing more robustly. Every private condition takes
$47.8$--$48.3$~s for the 100 rounds while the non-private baseline takes
$37.2$~s, a $1.2\%$ spread across private conditions against an $11$~s step up
from $\varepsilon=\infty$. Executed optimiser steps are identical across all
six conditions ($89$ per round), so this is not a workload difference.

The two mechanisms by which privacy consumes energy are therefore separable.
The per-round overhead is an implementation cost --- per-sample gradient
computation and clipping --- paid identically whether the budget is tight or
loose. Everything that depends on $\varepsilon$ acts through the number of
rounds, which the next result quantifies.

The engineering implications differ. A fixed overhead is addressed by
optimising the DP implementation; a convergence penalty is addressed by
changing the round budget or the stopping rule. Reporting only energy per
round would show privacy costing a flat $45\%$ and would miss the larger
effect entirely.

\begin{table}[t]
\caption{Energy per round on UCI-HAR, 100 rounds, means over three seeds}
\label{tab:perround}
\centering
\begin{tabular}{lrrrr}
\hline
$\varepsilon$ & Gross (W) & Net (W) & Gross (J) & Net (J) \\
\hline
0.5      & 62.30 & 44.09 & 29.84 & 21.12 \\
1        & 60.17 & 41.96 & 29.06 & 20.26 \\
2        & 60.79 & 42.58 & 29.36 & 20.56 \\
4        & 65.08 & 46.87 & 31.11 & 22.41 \\
8        & 62.26 & 44.05 & 30.07 & 21.27 \\
$\infty$ & 57.34 & 39.13 & 21.33 & 14.56 \\
\hline
\end{tabular}
\end{table}

\subsection{The Energy-Optimal Stopping Round}

This is the study's central result. Under a privacy budget the accuracy curve
peaks and then declines, and the peak arrives far earlier than the round
budget a non-private deployment would choose.

On HAR the non-private model needs the full budget to converge: it reaches
$0.886$ at round 95, having passed $95\%$ of that value only at round 59. At
$\varepsilon=1$ the same model on the same task peaks at $0.554$ at
\textbf{round 10} and then falls to $0.510$ by round 100. \textbf{The
energy-optimal stopping round is $9.5\times$ earlier under privacy}, and the
90 rounds that follow the peak --- $90\%$ of the run's energy --- buy a loss
of $4.3$ accuracy points.

The practical consequence is direct. A deployment that sizes its round budget
from the non-private convergence curve, which is the ordinary way to choose
one, spends roughly ten times the energy it needs under DP and ends with a
worse model than if it had stopped at round 10. The round budget must be
chosen from the privacy budget, not from the task.

The mechanism follows from the accounting model. Every round a client
participates in consumes budget, so $\sigma$ is set by the round budget before
training starts: on HAR at $n=279$ examples per client, $\sigma$ rises from
$10.16$ at 30 rounds to $18.44$ at 100. Extending training to reach
convergence is therefore self-defeating under local DP --- the extension is
what prevents convergence. Past the peak, the accumulated noise dominates the
aggregated update, and each further round subtracts signal.

% TODO: replace with the three-seed table from `analyze.py results_har/`.
% The two rows below are single-seed pilots (results_har_fix/), and the four
% remaining epsilon values are not yet filled in.
\begin{table}[t]
\caption{Energy-optimal stopping on UCI-HAR, 100 rounds. Rows marked $\dagger$
are single-seed pilots pending the sweep's own figures}
\label{tab:stopping}
\centering
\begin{tabular}{lrrrr}
\hline
$\varepsilon$ & Peak round & Peak acc. & Final acc. & Wasted (J) \\
\hline
0.5      & --- & --- & --- & --- \\
1$\dagger$        &  10 & 0.554 & 0.510 & --- \\
2        & --- & --- & --- & --- \\
4        & --- & --- & --- & --- \\
8        & --- & --- & --- & --- \\
$\infty\dagger$ &  95 & 0.886 & 0.880 & 0 \\
\hline
\end{tabular}
\end{table}

\subsection{Energy to Target Accuracy}

Energy per round makes privacy look cheap. What matters to a battery-powered
deployment is the energy needed to reach a usable model.

% TODO: fill from `analyze.py results_har/ --targets 0.40,0.50,0.85`.
% The default targets (0.50/0.70/0.85) leave every private condition unreached
% above 0.50, since eps=1 peaks at 0.554; lower targets give resolution inside
% the range the private conditions actually occupy, while 0.85 stays in the
% table to show the gap.

Unreachability is itself a result. A deployment that fixes both a privacy
budget and a round budget may be unable to produce a model of the required
quality at any energy expenditure. On HAR the non-private model reaches
$0.886$ while $\varepsilon=1$ never exceeds $0.554$: no additional training
closes that gap, because additional training is what raises $\sigma$. We
therefore report unreached targets explicitly rather than omitting those
configurations.

\begin{table}[t]
\caption{Energy to reach a target accuracy on UCI-HAR (J, net of idle).
``---'' marks targets not reached within 100 rounds}
\label{tab:target}
\centering
\begin{tabular}{lrrrr}
\hline
$\varepsilon$ & 0.40 & rounds & 0.50 & rounds \\
\hline
0.5      & --- & --- & --- & --- \\
1        & --- & --- & --- & --- \\
2        & --- & --- & --- & --- \\
4        & --- & --- & --- & --- \\
8        & --- & --- & --- & --- \\
$\infty$ & --- & --- & --- & --- \\
\hline
\end{tabular}
\end{table}

\subsection{CIFAR-10 as a Secondary Check}

CIFAR-10 over 30 rounds reproduces both effects on a harder task. The
per-round overhead is fixed at $+60\%$ net of idle ($121$~J against $75.8$~J),
flat in $\varepsilon$ to within $9.6\%$ against a $6.38\%$ floor. The peak
round again orders monotonically with the budget: 7, 11, 19, 24 and no peak
before round 30 at $\varepsilon \ge 8$. At $\varepsilon=0.5$ the model reaches
$0.169$ at round 7 and ends at $0.161$, with the following 23 rounds
accounting for $3866$~J, or $77\%$ of the run.

Reaching $0.20$ accuracy costs $1482$~J (net) at $\varepsilon=1$ against
$406$~J without privacy, $3.65\times$ for the same accuracy, because the
private run needs $11.7$ rounds where the non-private run needs $2.3$. At
$\varepsilon=0.5$ that target is never reached.

Absolute accuracies are low on this benchmark --- $0.503$ without privacy and
$0.16$--$0.31$ under it --- which is expected for a $62$k-parameter CNN over 30
non-IID rounds and is why HAR carries the headline claims. The agreement
between the two datasets on the shape of both effects is the point here.

\subsection{Ablations}

% TODO: the HAR ablation grid has not been run. The results below are
% CIFAR-10, three seeds per cell on the epochs and clipping axes.
% To add HAR: SEEDS="0 1 2" AXES="epochs clip" DATASET=har ROUNDS=100 ./run_ablations.sh

We vary one factor at a time at $\varepsilon \in \{1, \infty\}$ on CIFAR-10,
holding the others at the configuration of Section~III.

\textbf{Local epochs invert under privacy.} Increasing local work per round is
the standard way to reduce communication rounds in federated learning. Without
privacy it is close to accuracy-neutral here ($0.514$, $0.552$, $0.525$ for
one, two and five local epochs) while raising energy per round by $3.5\times$.
Under $\varepsilon=1$ it is harmful on both axes: accuracy falls from $0.181$
to $0.129$ to $0.101$, the peak round collapses from 13 to 6 to 3, and energy
per round rises from $164$~J to $585$~J. The accounting model explains this
directly: local steps enter $s_{\text{total}}$, so additional local epochs
consume budget faster and force a larger $\sigma$ for the same $\varepsilon$.
The two costs compound, since a configuration that spends more per round also
reaches its own peak sooner.

\textbf{Clipping norm trades accuracy against target.} At $\varepsilon=1$,
energy per round is unchanged by $C$ ($167$, $164$, $165$~J for $C \in \{0.5,
1, 2\}$, well within the noise floor), so $C$ acts purely on convergence.
Tighter clipping gives higher final accuracy ($0.206$ against $0.117$ at
$C=2$) and a later peak (round 18 against 6). The cheapest route to a
\emph{low} target is loose clipping --- $C=2$ reaches $0.20$ in $940$~J
against $1725$~J at $C=0.5$ --- but it then collapses, so that figure is only
meaningful together with the stopping round. The optimal clipping norm depends
on the accuracy target, and our fixed $C=1.0$ is a conservative choice rather
than an optimal one.

\textbf{Heterogeneity is masked by noise.} Without privacy, partition skew
behaves as expected: accuracy $0.397$, $0.509$, $0.547$ for $\alpha \in \{0.1,
0.5, 10\}$, and reaching $0.20$ costs $1.5\times$ more energy at $\alpha=0.1$
than at $\alpha=10$. Under $\varepsilon=1$ the ordering is not stable across
seeds, which we read as DP noise dominating the effect of heterogeneity rather
than as evidence of a direction. This axis does not arise on HAR, where the
partition is fixed by which volunteer contributed the data.

\subsection{Carbon Cost by Grid}

Identical workloads emit different amounts of CO$_2$ depending on where they
run. Table~\ref{tab:carbon} scales the measured energy to 1000 training runs
and applies published grid carbon intensities. The scaling repeats the
measured workload; it is not an extrapolation to different hardware.

% TODO: fill the HAR column set from `analyze.py results_har/`. The CIFAR-10
% figures below are final and may be kept as a second block or dropped.
On CIFAR-10 the cost attributable to privacy --- private minus non-private ---
is $228$~gCO$_2$eq per 1000 runs on Algeria's gas-dominated grid against
$12$~g in Norway, a factor of $18.8$ for the same computation and the same
guarantee. Choosing a privacy budget is therefore in part a geographic
decision: identical privacy requirements carry very different environmental
prices across deployment regions.

\begin{table}[t]
\caption{gCO$_2$eq per 1000 training runs on CIFAR-10. ``Privacy'' is the
private minus non-private difference}
\label{tab:carbon}
\centering
\begin{tabular}{lrrr}
\hline
Grid & DP & No DP & Privacy \\
\hline
Norway  &  36 &  24 &  12 \\
France  &  78 &  51 &  26 \\
Germany & 477 & 316 & 161 \\
USA     & 512 & 339 & 173 \\
Algeria & 677 & 448 & 228 \\
India   & 875 & 580 & 295 \\
\hline
\end{tabular}
\end{table}
